% !PRINTED 538
\weberpage{538}

We assume $a$ to be odd and let $\varphi$ run through a complete
system of representatives of the discriminant $-m$. If
then $m \equiv 3 \pmod 8$, then $c$ is odd, and $\varphi_1$, $\varphi_2$, $\varphi_3$
together run through a system of representatives of the discriminant $-4m$. For among the $3h$ roots of these forms
\[
\frac{\omega}{2}, \quad 2\,\omega, \quad 1 - \frac{2}{\omega+1}
\]
there occur, by § 123, no equivalent ones.

If on the other hand $m \equiv 7 \pmod 8$, then $c$ is even, $\frac{1}{2}\varphi_1$, $\frac{1}{2}\varphi_3$ each run through a system of representatives of the discriminant $-m$, while $\varphi_2$
runs through a system of representatives of the discriminant $-4m$.

If then $\psi$ runs through a system of representatives of the discriminant $-4m$, and we set
\weq{13}{S = \overset{x,y}{\Sigma}\, \frac{1}{\psi^s},}
then there results from (10):
\weq{14}{\begin{aligned}
\Bigl(1 + \frac{2}{4^s}\Bigr) \overset{\varphi}{\Sigma} S' &= \overset{\psi}{\Sigma} S, & m &\equiv 3 \\
\Bigl(1 + \frac{2}{4^s} - \frac{2}{2^s}\Bigr) \overset{\varphi}{\Sigma} S' &= \overset{\psi}{\Sigma} S, & m &\equiv 7
\end{aligned} \quad (\mathrm{mod}\ 8).}
If in formulae (29), (30) of the preceding section we set
\[
\psi' = 2\,\varphi, \quad \psi = (A, 2B, C),
\]
thus:
\[
\begin{aligned}
A &= a, \\
B &= b - a, \\
C &= 4c - 2b + a,
\end{aligned}
\]
then there results:
\[
S - \frac{1}{2^s} S' = \frac{2\pi}{\sqrt{m}} \log \frac{f(\omega_1)f(\omega_2)}{\sqrt{2}},
\]
and when $\psi$ runs through a system of representatives of the discriminant $-4m$,
then $\varphi$ runs three times, or only once, through a system of representatives of the discriminant $-m$, according as $m \equiv 3$
or $\equiv 7 \pmod 8$. We therefore obtain, if $\varepsilon$ in these
two cases is 3 or 1:
\weq{15}{\operatorname*{Lim}_{s=1} \Bigl(\overset{\psi}{\Sigma} S - \frac{\varepsilon}{2^s} \overset{\varphi}{\Sigma} S'\Bigr) = \frac{4\pi}{\sqrt{m}} \log \Pi \frac{f(\omega)}{\sqrt[4]{2}},}
in which $\omega$ runs through the roots of the forms $\psi$. And from this
there results, by (14):
%CONT>>
